\todo{Lecture 16}
\section{Extrema}
\begin{theorem}[Fermat]
Soient $I\subseteq \mathbb{R}$ un intervalle ouvert et $f:I\to \mathbb{R}$ derivable en $x_{0}\in I$. Si $f$ admet un extremum local en $x_{0}$, alors
$$
f'(x_{0})=0.
$$
\end{theorem}
\begin{proof}
Supposons que $x_{0}$ est un maximum local pour $f$. On peut choisir $\delta>0$ tel que :
\begin{itemize}
\item $(x_{0}-\delta,x_{0}+\delta)\subseteq I$ et
\item $|x-x_{0}|<\delta \Rightarrow f(x)\le f(x_{0})$.
\end{itemize}
Puisque $f$ est derivable en $x_{0}$, on a que
$$
f'(x_{0})=\lim_{ x \to x_{0} } \frac{f(x)-f(x_{0})}{x-x_{0}}.
$$
Posons $x_{n}=x_{0}-\frac{\delta}{n+1}$ et $y_{n}=x_{0}+\frac{\delta}{n+1}$, donc on a que $x_{n}\to x_{0}^{-}$ et $y_{n}\to x_{0}^{+}$. Donc
$$
f'(x_{0})=\lim_{ n\to \infty}  \frac{f(x_{n})-f(x_{0})}{x_{n}-x_{0}} \ge0
$$
et
$$
f'(x_{0})=\lim_{ n \to \infty } \frac{f(y_{n})-f(x_{0})}{y_{n}-x_{0}} \le 0.
$$
Donc $f'(x_{0})=0$.
\end{proof}
\begin{theorem}[Rolle]
Soient $a<b$ et $f:[a,b]\to \mathbb{R}$. Si
\begin{itemize}
\item $f$ est continue sur $[a,b]$;
\item $f$ est derivable sur $(a,b)$ et
\item $f(a)=f(b)$,
\end{itemize}
alors il existe $c\in(a,b)$ tel que $f'(c)=0$.
\end{theorem}
\begin{proof}
Par le theoreme de Weierstrass, $f$ atteint son minimum et son maximum sur $[a,b]$. Posons
$$
m=\min_{x \in [a,b]}f(x),\quad M=\max_{x \in [a,b]} f(x).
$$
Si $m=M$, alors $f$ est constante sur $[a,b]$, donc $f'(c)=0$ pour tout $c\in(a,b)$. Sinon, vu que $f(a)=f(b)$, on a que $m$ ou $M$ est atteint en un point $c\in(a,b)$ et donc par Fermat on conclut. 
\end{proof}
\begin{theorem}[Theoreme des accroisements tinis : TAF]
Soient $a<b$ et $f:[a,b]\to \mathbb{R}$. Si $f$ est continue sur $[a,b]$ et derivable sur $(a,b)$, alors il existe $c\in(a,b)$ tel que
$$
f'(c)= \frac{f(b)-f(a)}{b - a}. 
$$
\end{theorem}
\begin{proof}
On soustrait a $f$ l'equation de la droite qui passe par $(a,f(a))$ et $(b,f(b))$ :
$$
g(x):= f(x)-f(a)-\frac{f(b)-f(a)}{b - a}(x-a).
$$
On a $g(a)=g(b)=0$ et $g$ est continue sur $[a,b]$ et derivable sur $(a,b)$. On a donc par Rolle, qu'il existe $c\in(a,b)$ tel que $g'(c)=0$. Donc
$$
0=g'(c)=f'(c)-\frac{f(b)-f(a)}{b - a}.
$$
\end{proof}

\begin{definition}
L'interieur d'un intervalle $I$ est l'intervalle $\mathrm{int}(I)=I\setminus \{ \inf I, \sup I \}$. Par exemple, $\mathrm{int}([a,b])=(a,b)$, $\mathrm{int}((a,b))=(a,b)$ et $\mathrm{int}(\mathbb{R})=\mathbb{R}$.
\end{definition}
\begin{corollary}
Soient $I$ un intervalle non-reduit a un point et $f:I\to \mathbb{R}$ continue sur $I$ et derivable sur son interieur $\mathrm{int}(I)$. Alors, $f$ est :
\begin{itemize}
\item croissante sur $I$ si et seulement si $f'(x)=0$ pour tout $x \in \mathrm{int}(I)$;
\item croissante sur $I$ si et seulement si $f'(x)\ge 0$ pour tout $x \in \mathrm{int}(I)$;
\item decroissante sur $I$ si et seulement si $f'(x)\le 0$ pour tout $x \in \mathrm{int}(I)$;
\item strictement croissante sur $I$ si $f'(x)> 0$ pour tout $x \in \mathrm{int}(I)$.
\end{itemize}
\end{corollary}
\begin{proposition}
Soient $I$ un intervalle ouvert et $f:I\to \mathbb{R}$ une fonction derivable. Supposons que $f'$ est derivable en $x_{0}\in I$.
\begin{itemize}
\item Si $f'(x_{0})=0$ et $f''(x_{0})>0$, alors $x_{0}$ est un minimum local.
\item Si $f'(x_{0})=0$ et $f''(x_{0})<0$, alors $x_{0}$ est un maximum local.
\end{itemize}
\end{proposition}
\begin{proof}
Supposons que $f'(x_{0})=0$ et $f''(x_{0})>0$. Il faut montrer que $f$ est decroissante a gauche de $x_{0}$ et croissante a droite de $x_{0}$. Par definition
$$
0<f''(x_{0})=\lim_{ x \to x_{0} }  \frac{f'(x)-f'(x_{0})}{x-x_{0}}.
$$
Donc, il existe $\delta>0$ tel que $(x_{0}-\delta,x_{0}+\delta)\subset I$ et 
$$
\frac{f'(x)}{x-x_{0}}> \frac{f''(x_{0})}{2}>0,\quad \forall x:0<|x-x_{0}|<\delta.
$$
Pour $x \in(x_{0},x_{0}+\delta)$, cela implique que $f'(x)>0$, donc $f$ est croissante sur $(x_{0},x_{0}+\delta)$. Pour $x \in(x_{0}-\delta,x_{0})$, on a $f'(x)<0$, donc $f$ est décroissante sur ce intervalle.
\end{proof}
\begin{theorem}
Soit $f:I\to \mathbb{R}$ une fonction derivable sur un intervalle $I$, et soit $L\ge 0$. Les deux proprietes suivantes sont equivalentes :
\begin{enumerate}
\item $f$ est $L$-lipschitzienne sur $I$;
\item $|f'(x)|\le L$ pour tout $x \in I$.
\end{enumerate}
\end{theorem}
\begin{proof}
\begin{itemize}
\item $2\Rightarrow 1$ : Soient $x,y\in I$ avec $x\neq y$. Par le TAF, il existe $c\in(x,y)$ tel que 
$$
f(x)-f(y)=f'(c)(x-y)
$$
et donc
$$
|f(x)-f(y)|=|f'(c)||x-y|\le L|x-y|.
$$
\item $1\Rightarrow 2$ : Soit $x_{0}\in I$. On a
$$
\left\lvert  \frac{f(x)-f(x_{0})}{x-x_{0}}  \right\rvert \le \frac{L|x-x_{0}|}{|x-x_{0}|}=L.
$$
Donc $|f'(x_{0})|\le L$.
\end{itemize}
\end{proof}